\documentclass[10pt,a4paper]{amsart}
\usepackage[margin=19mm]{geometry}
\usepackage{amsmath,amssymb,mathtools}
\usepackage[scr=rsfs,cal=euler]{mathalfa}
\usepackage{xfrac}
\usepackage[unicode,hidelinks,bookmarks=false]{hyperref}
\newcommand{\ie}{i.e.\ }
\newcommand{\cf}{cf.\ }
\newcommand{\resp}{respectively\ }
\newcommand{\Subsection}{\subsection}
\providecommand{\nobreakdash}{\nobreak\hbox{-}}
\setlength{\emergencystretch}{2em}
\multlinegap0pt
\usepackage{bm}
\usepackage{bbm}
\usepackage{dsfont}
\usepackage{xspace}
\usepackage{cancel}
\usepackage{mathtools}
\usepackage{amsthm}
\makeatletter
\@ifundefined{theorem}{\newtheorem{theorem}{Theorem}}{}
\@ifundefined{proposition}{\newtheorem{proposition}[theorem]{Proposition}}{}
\makeatother
\title[An Effective Method for Solving a Class of Transcendental Di]{An Effective Method for Solving a Class of Transcendental Diophantine Equations}
\author{Zeyu Cai}
\date{Source: arXiv:2510.11753v1 (CC BY 4.0)}
\begin{document}
\maketitle

\noindent\emph{Source and licence.} An Effective Method for Solving a Class of Transcendental Diophantine Equations. Version arXiv:2510.11753v1 (CC BY 4.0), distributed under the Creative Commons Attribution 4.0 International licence, as recorded in the arXiv licence metadata for this exact version and shown on the abstract page https://arxiv.org/abs/2510.11753v1. Full rights notice: licenses/arxiv-2510.11753-CC-BY-4.0.txt.\par\smallskip

\noindent\textit{Editorial scope.} Counted unit: the Proposition whose source label is \texttt{None} in the article (source0.tex lines 64--66, source label \texttt{None}), delivered together with the article's own complete original proof of it (source0.tex lines 67--82, one \texttt{proof} environment of 16 lines). No mathematics is written, repaired or re-derived: statement and proof are the article's own text line for line. Required context retained uncounted: none. Every definition, display and label of those lines is retained unchanged. Omitted from this package and not redistributed: 1--63, 83--900. Nothing inside a retained excerpt is omitted or altered: every retained body is delivered exactly as the article's lines are written, so no omission receipt of any kind accompanies this package. Citations: the retained excerpts carry no citation command at all, so no bibliography entry of the article is transcribed and no reference is redistributed. Reference adaptation: none was needed and none was made; every reference inside the delivered unit resolves inside it, so no unresolved reference or citation is produced. Printed slips kept literal: no printed word, symbol, hypothesis or number of the counted statement or of its proof was altered. Layout and class: the article's own class is \texttt{article} with packages including geometry, amsmath, amssymb, amsfonts, amsthm, graphicx, booktabs, natbib, hyperref, fontenc, minted; the wrapper is the lane's pinned \texttt{amsart} editorial wrapper at 10pt on A4, which restates the theorem-like environments the retained text needs and adds the article's own macro definitions that the retained text needs (none), with extra wrapper packages (bm, bbm, dsfont, xspace, cancel, mathtools). The delivered numbering is the unit's own. Assets: no retained span contains a figure environment, a graphics-inclusion command, a PDF-page inclusion, a table or a TikZ picture; no image file is copied into this package, the package declares no asset dependency and no image file is redistributed with it. Licence: version arXiv:2510.11753v1 of this article is distributed under the Creative Commons Attribution 4.0 International licence, as recorded in the arXiv licence metadata for this exact version and shown on the abstract page https://arxiv.org/abs/2510.11753v1. A full rights notice is delivered at licenses/arxiv-2510.11753-CC-BY-4.0.txt. Characters: every retained excerpt is pure ASCII, so the delivered engine, XeLaTeX with its default Latin Modern text font, reports no missing character.
\begin{proposition}
\[ 5^x+3=2^y, \quad x, y \in \mathbb{N}^* \Rightarrow (x,y)=(1,3) \text{ or } (x,y)=(3,7) \]
\end{proposition}

\begin{proof}
By checking small positive integer solutions, we confirm that $(1, 3)$ and $(3, 7)$ are two solutions to the equation. To prove that no other solutions exist, we assume there is a solution $(x, y)$ satisfying $y \ge 8$.
Consider the equation modulo $2^8=256$. When $y \ge 8$, $2^y \equiv 0 \pmod{256}$. The original equation becomes:
$$ 5^x+3 \equiv 0 \pmod{256} \implies 5^x \equiv 253 \pmod{256} $$
Solving this exponential congruence equation, we get $x \equiv 35 \pmod{64}$. This means $x$ must satisfy one of the following four cases: $x \equiv 35 \pmod{256}$, $x \equiv 99 \pmod{256}$, $x \equiv 163 \pmod{256}$, or $x \equiv 227 \pmod{256}$.

Next, consider the equation modulo $257$ (a prime number).
\begin{itemize}
    \item If $x \equiv 35 \pmod{256}$, then $5^x \equiv 14 \pmod{257}$. The original equation gives $2^y = 5^x+3 \equiv 17 \pmod{257}$.
    \item If $x \equiv 99 \pmod{256}$, then $5^x \equiv 224 \pmod{257}$. The original equation gives $2^y = 5^x+3 \equiv 227 \pmod{257}$.
    \item If $x \equiv 163 \pmod{256}$, then $5^x \equiv 243 \pmod{257}$. The original equation gives $2^y = 5^x+3 \equiv 246 \pmod{257}$.
    \item If $x \equiv 227 \pmod{256}$, then $5^x \equiv 33 \pmod{257}$. The original equation gives $2^y = 5^x+3 \equiv 36 \pmod{257}$.
\end{itemize}
However, by computing the powers of $2$ modulo $257$, it can be verified that none of $17, 227, 246, 36$ are in the set $\{2^k \pmod{257} | k \in \mathbb{N}^*\}$. This is a contradiction.
Therefore, the assumption $y \ge 8$ is false. By enumerating $y \in \{1, 2, \dots, 7\}$, we can verify that there are no other solutions.
\end{proof}

\end{document}
